% Vietnamese translation for OI Public Library
% Source-File: IOI中国国家候选队论文/2009/徐持衡/浅谈几类背包题.ppt
\documentclass[11pt]{article}
\usepackage{oipl-vn}

\newcommand{\Oh}{\mathcal{O}}

\title{Bàn sơ lược về vài dạng bài ba lô\\{\large Ghi chú theo slide}}
\author{Xu Chiheng\\Trường Trung học Ôn Châu, Chiết Giang\\Người hướng dẫn: Shu Chunping}
\date{Tháng 12 năm 2008}

\begin{document}
\maketitle

\origtitle{Several variants of knapsack problems}
\sourcepath{IOI-China-Candidate-Team-Papers/2009/Xu-Chiheng/knapsack-slides.ppt}

\section{Mạch trình bày}

Slide bắt đầu từ ba lô \(0/1\), ba lô hoàn toàn và ba lô nhiều lần, sau đó
đưa người đọc đến các biến thể có quan hệ phụ thuộc hoặc trạng thái khó hơn.
Trục chính là: viết đúng công thức chuyển, rồi giảm số lần chuyển bằng cấu
trúc dữ liệu hoặc đổi cách nhóm vật.

\section{Các kỹ thuật chính}

\begin{itemize}
  \item Ba lô hoàn toàn dùng thứ tự duyệt tăng theo dung lượng để cho phép
  dùng lại cùng một vật nhiều lần.
  \item Ba lô nhiều lần có thể tách nhị phân số lượng, biến thành một số vật
  \(0/1\) với tổng kích thước logarit.
  \item Khi công thức có dạng lấy cực trị trên một cửa sổ theo lớp đồng dư,
  hàng đợi đơn điệu giảm độ phức tạp của mỗi nhóm từ tuyến tính nhân số lượng
  xuống tuyến tính theo dung lượng.
  \item Với ba lô phụ thuộc dạng cây, trạng thái của nút cha được ghép từ các
  con bằng quy hoạch động trên cây.
\end{itemize}

\section{Cách đọc ví dụ}

Các ví dụ trong bài viết đầy đủ cung cấp nhiều chi tiết đề bài. Ở bản slide,
nên tập trung vào dấu hiệu nhận dạng: vật có thể chọn bao nhiêu lần, có quan
hệ phụ thuộc nào, và phần cực trị trong chuyển trạng thái có thể được duy trì
bằng một cấu trúc đơn điệu hay không.

\end{document}
